\chapter{Performance Tuning}
\label{ch:performance}

\newthought{The runtime is sized for the slow regime}: external calculators taking
seconds to minutes, results landing on network storage. In that regime the framework's
overhead is noise, and performance is a resource-allocation problem with three dials ---
\yamlkey{workers}, \yamlkey{Pools}, and the archiving policy. This chapter is the tuning
guide, ordered by how often each dial actually matters.

\section{First, measure}
\label{sec:perf-measure}

Two numbers from the \code{[Scan Performance]} block
(Section~\ref{sec:scan-performance}) frame everything:

\begin{quote}
\emph{effective parallelism} $\;=\;$
\code{avg\_point\_eval\_sec} $/$ \code{avg\_sample\_sec}
\end{quote}

If a point costs $60$\,s to evaluate and the system delivers one finished sample every
$2$\,s, you are running $\sim$$30$ points wide. Compare that against
\yamlkey{workers}: a large gap means Workers are waiting --- on pool slots, on the
queue, or on storage --- and the live monitor shows which
(Section~\ref{sec:monitor-diagnosis}).

\section{Dial 1: \yamlkey{workers}}
\label{sec:perf-workers}

Workers are cheap (a Python process holding calculator templates); their \emph{work} is
not. Size by the binding resource, not by optimism:

\begin{itemize}
  \item \textbf{CPU-bound calculators}: about one Worker per available core (times the
        cores each calculator uses).
  \item \textbf{Memory-bound}: total \textsc{ram} $/$ peak per-point footprint ---
        calculators, not Workers, dominate the footprint.
  \item \textbf{Subprocess-wait-bound} (calculators that themselves farm out, or wait
        on I/O): oversubscribe cores freely; the \textsc{gil} is released during waits.
\end{itemize}

Scaling is horizontal and honest: doubling Workers doubles throughput until a pool,
the machine, or storage saturates. The design target is hundreds of Workers.

\begin{jnote}
More Workers also means more one-time initialization (template loads, environment
captures, expression compiles) and --- for \yamlkey{clone\_shadow} tools --- more pack
installs on \emph{first contact} with each pack. Long scans amortize all of it; a
five-minute scan with 128 Workers spends its life starting up.
\end{jnote}

\section{Dial 2: calculator \yamlkey{Pools}}
\label{sec:perf-pools}

Pools decouple ``how many Workers exist'' from ``how many copies of tool X may run''
(Section~\ref{sec:pools}). Patterns:

\begin{itemize}
  \item \textbf{The license/memory hog}: pool $=$ what you can afford; other layers'
        work continues around it.
  \item \textbf{The mismatched chain}: a 10\,s tool feeding a 10\,min tool wants pools
        sized $\sim$1:60 against Worker count, not equal shares.
  \item \textbf{Diagnosis}: monitor shows every slot of one pool busy while Workers
        idle $\Rightarrow$ that pool is the throttle. Raise it or accept it as the
        design capacity.
\end{itemize}

A Worker that waits longer than the internal 30\,s slot timeout fails that Sample ---
chronically starved pools show up as timeout failures, not just slowness.

\section{Dial 3: archiving and storage}
\label{sec:perf-storage}

Storage is the usual silent bottleneck on clusters:

\begin{itemize}
  \item \textbf{Watch the archive queue} in the monitor: steady growth means the
        Archiver cannot keep up with the Workers.
  \item \yamlkey{archiver.batch\_size} (default 200) trades database write frequency
        against memory; larger batches help slow filesystems.
  \item \textbf{Buckets} (default 200/dir, packed) exist precisely for network
        filesystems --- do not disable \yamlkey{pack} on \textsc{nas} scans that
        \yamlkey{save} files.
  \item \textbf{Point the output root at the fast filesystem} rather than trying to
        buffer through one. Workers write each Sample's products straight to their final
        location under \file{SAMPLE/} (Section~\ref{sec:sample-lifecycle}), so a remote
        or slow output tree is paid for on the critical path and no task-card setting
        moves that hop elsewhere. Run against local scratch and copy the finished output
        directory afterwards.
  \item \textbf{Save less} (Section~\ref{sec:save-policy}): the cheapest byte is the
        one never written.
\end{itemize}

\section{What not to bother tuning}
\label{sec:perf-non-dials}

\begin{itemize}
  \item \yamlkey{batch\_size} (submit grouping): queue pipelining only; the default
        256 is fine from ten samples to ten million.
  \item \textbf{Expression and template caching}: automatic, once per Worker.
  \item \textbf{Redis}: at HEP scan rates the broker idles; it is never the bottleneck
        on a single machine.
  \item \textbf{Opera micro-costs}: an opera call is microseconds against calculator
        seconds. If your workflow is \emph{all} operas and you need raw throughput,
        you are outside the slow regime this runtime is sized for --- it will still do
        thousands of samples per minute, but that is the ceiling's neighborhood.
\end{itemize}

\begin{jtip}
The tuning loop that works: run 15 minutes at a plausible \yamlkey{workers}; read
\code{[Scan Performance]} and one monitor snapshot; move \emph{one} dial; repeat. Two
iterations usually land within 20\% of the machine's ceiling --- past that, buy
hardware, not \textsc{yaml}.
\end{jtip}
